\scanpage{7}
Mesures spectrales et réarrangements de fonctions.

\margin{\(\unsure{l}(M),\ L^1(M),\ L^\infty(M)\)\par \(K(M)\)\par \(\unsure{K^\infty(M)|K^\alpha(M)}\)\par \(C_0(M)\)\par \deleted{\(\unsure{\Phi|R}\)}}

Soit \(f\) une fonction mesurable complexe sur l'espace mesuré \(M\) (muni de la mesure \(m\)). On dit que \(f\) \uline{admet une mesure spectrale} si l'image de \(m\) par \(f\) en tant que mesure sur \(\underline C^*\) existe. Cela signifie donc aussi que pour \(0<a<b<+\infty\), \(|f|^{-1}([a,b])\) est intégrable (donc cette condition ne dépend que de \(|f|\)). \(f(m)\) (considérée comme mesure sur \(C^*\)) est appelée \uline{mesure spectrale de \(f\)}, et notée \(\mu_f\) :
\[
(1)\qquad \mu_f=f(m)\qquad\text{i.e.}
\]
\[
(2)\qquad\langle\alpha,\mu_f\rangle=\langle\alpha\circ f,m\rangle=\int\alpha\circ f\,dm
\qquad\deleted{\unsure{\alpha\circ f\,dm}}
\]
pour toute fonction continue à support \unsure{compact} sur \(\underline C^*\). \(\mu_f\) est \uline{positive}. Elle est portée par \(\underline R_*\) (resp. par \(R_*^+\)) \uline{si et seulement si \(f\) est réelle} (resp. positive). \deleted{\unsure{Deux}} Une fonction \(f\) sur l'espace mesuré \(M\) et une fonction \(g\) sur l'espace mesuré \(N\) sont dites \uline{spectralement équivalentes} si elles admettent une même mesure spectrale.

\uuline{Exemples} Si l'ensemble \(f^{-1}(\unsure{\{\ne0\}|\complement\{0\}})\) est intégrable, alors \(\mu_f\) existe et est bornée. C'est en particulier le cas si \(f\) est à support compact, ou si \(m\) est bornée.
\inserted{\unsure{Autre cas : toute} \(f\in C_0(M)\), \unsure{plus généralement telle} que, \deleted{\(\unsure{\varepsilon}\)} pour tout \(\varepsilon\), \(|f|^{-1}(]\varepsilon,\infty[)\) est intégrable.}

\margin{\unsure{continuité}\par\unsure{à gauche}}
\uline{Proposition 1} \deleted{\(\unsure{0}\)} Une fonction positive décroissante \(\varphi\) sur \(R_*^+\) \(\bigl[\) a une mesure spectrale si et seulement si elle est nulle à l'infini; \(\mu_\varphi\) est une mesure positive sur \(R_*^+\) telle que
\[
\int_1^\infty d\mu_\varphi<+\infty.
\]
Réciproquement, pour toute mesure positive \(\nu\) sur \(R_*^+\) telle que \(\int_1^\infty d\nu<+\infty\), existe une fonction \inserted{positive} décroissante \(\varphi\) et une seule (à une \uline{modification} \deleted{\unsure{sur un ensemble}}\inserted{négligeable} près) \unsure{qui admette} \(\nu\) comme mesure spectrale.

On a donc pour presque tout \(t\)
\[
(3)\qquad t=\int_s^\infty d\unsure{\mu|\nu}(s)\qquad(\text{si }s=\varphi(t))
\]

\mspage{8}
